\documentclass[a4paper]{article}

\usepackage[fontset=fandol]{ctex}
\usepackage{amsmath, amssymb}
\usepackage{graphicx}
\graphicspath{{.}{../}}
\usepackage[margin=2.5cm]{geometry}
\usepackage[most]{tcolorbox}
\usepackage{etoolbox}
\usepackage{listings}
\usepackage{hyperref}
\usepackage{booktabs}
\usepackage{subcaption}
\usepackage{float}
\usepackage{tikz}
\IfFileExists{pgfplots.sty}{
    \usepackage{pgfplots}
    \pgfplotsset{compat=1.18}
}{}

\newtcolorbox{knowledgebox}[1]{
    enhanced, colback=blue!5!white, colframe=blue!75!black, colbacktitle=blue!75!black,
    coltitle=white, fonttitle=\bfseries, title=#1, attach boxed title to top left={yshift=-2mm, xshift=2mm},
    boxrule=1pt, sharp corners
}

\newtcolorbox{importantbox}[1]{
    enhanced, colback=yellow!10!white, colframe=yellow!80!black, colbacktitle=yellow!80!black,
    coltitle=black, fonttitle=\bfseries, title=#1, sharp corners
}

\newtcolorbox{warningbox}[1]{
    enhanced, colback=red!5!white, colframe=red!75!black, colbacktitle=red!75!black,
    coltitle=white, fonttitle=\bfseries, title=#1, sharp corners
}

\lstset{
    language=Python,
    basicstyle=\ttfamily\small,
    keywordstyle=\color{blue},
    stringstyle=\color{red!60!black},
    commentstyle=\color{green!60!black},
    breaklines=true,
    frame=single,
    numbers=left,
    numberstyle=\tiny\color{gray},
    captionpos=b,
    extendedchars=false
}

\newcommand{\notetitle}{课程笔记：Lecture 6B 强化学习（Reinforcement Learning）}
\newcommand{\noteauthors}{Codex}
\newcommand{\notedate}{\today}
\newcommand{\videochannel}{Robots That Learn / UC Berkeley CS 294-277}
\newcommand{\videopublishdate}{课程讲次日期：2024-10-14}
\newcommand{\videoduration}{39:30}
\newcommand{\videourl}{https://www.youtube.com/watch?v=pvTYmRXLixo}
\newcommand{\videocoverpath}{cover.jpg}

\begin{document}

\begin{titlepage}
\centering
{\Large 课程笔记\par}
\vspace{1.2cm}
{\huge\bfseries \notetitle\par}
\vspace{0.8cm}
{\large \noteauthors\par}
\vspace{0.3cm}
{\large \notedate\par}
\vspace{1.2cm}

\ifdefempty{\videocoverpath}{
{\small 请在生成笔记时填入视频封面图的本地路径。\par}
}{
\includegraphics[width=0.82\textwidth,height=0.45\textheight,keepaspectratio]{\videocoverpath}\par
}

\vfill
\begin{tcolorbox}[width=0.9\textwidth, colback=black!2!white, colframe=black!60, sharp corners]
\textbf{视频作者/频道}：\videochannel\par
\textbf{发布日期}：\videopublishdate\par
\textbf{视频时长}：\videoduration\par
\textbf{视频链接}：\href{\videourl}{\nolinkurl{\videourl}}
\end{tcolorbox}
\end{titlepage}

\tableofcontents
\newpage

\begin{figure}[H]
\centering
\includegraphics[width=0.78\textwidth]{overview_cover.jpg}
\caption{课程官方材料中的代表性页面，用于帮助读者在进入本讲正文前建立整体上下文。}
\end{figure}
\section{课程定位与本讲目标}

这一讲延续上一讲的 RL 主线，但讨论重心明显从“算法是什么”移动到“怎样把算法真正用到机器人里”。讲者想回答的核心问题是：如果我们已经知道 RL 可以学策略，为什么机器人领域仍然如此依赖模拟器、并行仿真和 domain randomization（域随机化）？

\subsection{本讲在课程中的位置}

上一讲讲的是 RL 的基本定义、policy gradient 族和 simulator 的物理骨架；这一讲则进一步进入实践层：有哪些 simulators 被广泛使用，它们各自适合什么任务，sim-to-real gap 到底如何缩小，以及现实世界里有哪些代表性成功案例。

\begin{importantbox}{本讲的关键判断}
在机器人 RL 里，“训练快”不是唯一标准。精度、可控性、可并行性、接口成熟度和部署便利性，通常比单纯的仿真速度更重要。
\end{importantbox}

\subsection{本章小结}

如果把上一讲看成“RL 的算法底座”，那么这一讲就是“RL 的工程底座”。它关心的不只是模型能否学会，而是模型能否跨过仿真与现实之间的鸿沟。

\section{Simulation Environments：哪些模拟器在用，为什么}

\subsection{MuJoCo、Raisim、Brax 与 Isaac Gym}

课程列举了几种学术界和工业界常见的 simulator：
\begin{itemize}
  \item MuJoCo（Multi-Joint dynamics with Contact）：广泛使用的物理引擎，擅长多刚体与接触建模，背后以 C++ 实现，近年来也有面向 GPU/TPU 并行的 MuJoCo XLA（MJX）版本；
  \item Raisim：偏 CPU 的刚体物理引擎，适合机器人与 AI 研究；
  \item Brax：用 JAX 写成、面向加速硬件的可微物理引擎，适合单设备和多设备并行；
  \item Isaac Gym：NVIDIA 推出的 GPU-based simulator，强调大规模并行仿真。
\end{itemize}

这些选择背后的共同点是：它们都试图让机器人训练更快、更标准化、更可重复，但它们侧重点不同。MuJoCo 偏通用且成熟，Brax 偏研究与可微分，Isaac Gym 偏并行和 GPU 吞吐，而 Raisim 则常见于 CPU 侧的高精度刚体场景。

\subsection{速度不等于最好}

讲者特别提醒，simulator 不是越快越好。某些任务需要更高的数值精度，或者对接触建模更敏感，此时 CPU 方案反而可能比 GPU 更合适。换句话说，simulator 的选择应该由任务决定，而不是由“谁跑得最快”单独决定。

\begin{warningbox}{不要把 simulator 只当成加速器}
很多失败的机器人实验不是因为算法错了，而是因为模拟器与任务不匹配：接触模型、数值稳定性、精度和接口方式都会影响最终的策略质量。
\end{warningbox}

\subsection{本章小结}

模拟器的价值不只在吞吐量，更在于它是否能提供足够可信、足够稳定、足够可控的训练环境。对机器人 RL 来说，训练平台本身就是方法的一部分。

\section{Sim2Real 的两大工具：Parallel Simulation 与 Domain Randomization}

\subsection{Parallel Simulation}

Parallel simulation 的作用是一次生成大量交互数据。课程强调，这一点对 RL 尤其关键，因为 RL 的学习通常依赖大量探索，而单线程、单环境采样会让训练速度非常慢。并行仿真可以把多个环境同时跑起来，从而显著提高样本生成速度。

讲者给出的直觉是：如果一个环境只能顺序地产生经验，那么策略改进会被物理世界的时间约束卡住；如果能把训练搬到大量并行环境里，就可以把“等机器人动起来”的时间转化成“算力在同一时刻处理更多轨迹”。

\subsection{Domain Randomization}

Domain randomization 的核心是：在训练时故意打乱模拟环境的若干方面，让策略见过更广的变化，从而在真实世界里更容易泛化。课程把它分成两类：
\begin{itemize}
  \item dynamics randomization：随机化动力学参数，如质量、摩擦、外力扰动；
  \item visual randomization：随机化视觉因素，如纹理、光照、相机位置、噪声和模糊。
\end{itemize}

这一思想本质上是在做“把现实世界的不确定性预先摊开”。如果真实机器人面对的参数会漂移，那么训练时就不要假设参数固定不变。

\begin{knowledgebox}{域随机化的直觉}
不是让仿真更像某一个现实，而是让策略对一大类现实都“见怪不怪”。只要真实系统落在训练分布附近，策略就更可能稳定工作。
\end{knowledgebox}

\subsection{本章小结}

Parallel simulation 解决“数据生成太慢”的问题，domain randomization 解决“模拟过于理想化”的问题。它们通常一起出现，因为一个负责规模，一个负责鲁棒性。

\section{案例一：OpenAI dexterous in-hand manipulation}

\subsection{为什么这个案例重要}

课程把 OpenAI 的灵巧手操作案例当成 sim-to-real 的代表性成功之一。原因很简单：dexterous in-hand manipulation（灵巧手内操作）是一个接触复杂、动力学高度耦合、现实误差很难忽略的任务。如果连这种任务都能从仿真迁移到现实，那么 sim-to-real 的方法论就有了很强说服力。

\subsection{这个案例里怎么做 domain randomization}

讲者列出的随机化内容非常具体：
\begin{itemize}
  \item observation noise：给观测加入高斯噪声，同时加入相关噪声和非相关噪声；
  \item physics randomizations：对质量、摩擦、关节阻尼等参数进行采样；
  \item timing noise：通过改变每一步对应的物理仿真步数，模拟执行节奏变化；
  \item unmodeled effects：把仿真中没有显式建模的效果补进去，例如 tendon actuation 带来的 backlash（回程间隙）；
  \item visual appearance randomizations：改变灯光、相机、纹理和材料外观。
\end{itemize}

这里最值得注意的是，很多随机化并不是“加点噪声”这么简单，而是针对真实系统中最容易失配的环节逐项补偿。比如 backlash 不是视觉问题，而是机械传动的物理问题。

\begin{importantbox}{真实部署时最难的地方}
灵巧操作的难点不在“会不会做动作”，而在接触、误差、时间尺度和传感器噪声同时出现时，策略还能不能保持稳定。
\end{importantbox}

\subsection{本章小结}

OpenAI 灵巧手案例说明，sim-to-real 并不是把一个仿真策略直接扔到现实里，而是要通过系统性的随机化，把仿真中的假设尽量推近现实中的不确定性分布。

\section{案例二：High-speed quadrotor flight in the wild}

\subsection{为什么这个案例也能成立}

第二个案例是高速无人机在真实环境中的飞行。这个问题的特点是：决策必须快、视觉输入可能复杂、环境变化大，而且对在线推理延迟极其敏感。课程用它说明，sim-to-real 的关键不只在“把策略训出来”，还在“让感知与控制的接口足够轻量”。

\subsection{四个设计选择}

课程把这个案例的成功归纳为四点：
\begin{enumerate}
  \item 使用 depth image（深度图）而不是 RGB 图像作为输入，降低视觉外观差异；
  \item 采用轻量级神经网络，减少机载推理延迟；
  \item 让网络输出多模态轨迹和碰撞代价，而不是单一动作；
  \item 训练完成后 zero-shot 地直接部署到真实环境。
\end{enumerate}

其中 depth representation 很关键，因为它对纹理、颜色和光照变化不那么敏感，更贴近“几何结构”而不是“外观”本身。多模态输出也很重要，因为高速飞行中的可行路径往往不止一条，强行平均反而会得到危险的中间解。

\subsection{为什么讲者说 RL 不是唯一选项}

这部分课件和讲解的一个微妙但重要的信息是：能成功并不意味着必须靠纯 RL。对于某些高速导航任务，感知设计、轨迹预测和损失函数设计同样关键。讲者强调这一点，是为了避免把 RL 神化成万能方案。

\subsection{本章小结}

无人机案例说明，sim-to-real 的成功常常来自一组组合拳：合适的输入表示、轻量模型、多模态输出、合理训练和足够鲁棒的部署策略，而不是单一技巧。

\section{总结与延伸}

这一讲可以总结为一句非常工程化的话：机器人 RL 想要真正落地，必须同时处理“算得快”和“迁得过去”这两件事。前者依赖并行仿真和合适的 simulator，后者依赖 domain randomization、输入表示设计和任务级的鲁棒性策略。

\subsection{值得记住的三点}

\begin{itemize}
  \item 模拟器的选择取决于任务，而不是单纯取决于速度；
  \item parallel simulation 解决吞吐，domain randomization 解决失配；
  \item 高质量 sim-to-real 往往是系统设计，而不只是算法技巧。
\end{itemize}

\subsection{拓展阅读}

\begin{itemize}
  \item \emph{Learning Dexterous In-Hand Manipulation}
  \item \emph{Learning agile and dynamic motor skills for legged robots}
  \item \emph{Learning High-Speed Flight in the Wild}
\end{itemize}

\subsection{结束语}

如果把上一讲理解为“RL 作为算法”，这一讲就更像“RL 作为系统”。真正强的机器人系统通常不是只靠一个模型，而是靠 simulator、randomization、表示学习和部署约束共同组成的整体。

\end{document}
